% =====================================================================
%  TÀI LIỆU KIỂM THỬ PHẦN MỀM — HỆ THỐNG NAVIMART
%  Biên dịch: XeLaTeX (khuyến nghị) hoặc LuaLaTeX để hỗ trợ tiếng Việt UTF-8.
%      xelatex TestingReport-NaviMart.tex   (chạy 2 lần để cập nhật mục lục)
%  Nếu thiếu dấu tiếng Việt, mở khối FONT bên dưới và chọn font có sẵn.
% =====================================================================
\documentclass[12pt,a4paper]{article}

\usepackage{fontspec}
% ---- FONT (tùy chọn) -------------------------------------------------
% \setmainfont{Times New Roman}
% \setmonofont{Consolas}
% ---------------------------------------------------------------------

\usepackage[a4paper,margin=2.5cm]{geometry}
\usepackage{tabularx}
\usepackage{longtable}
\usepackage{array}
\usepackage[table]{xcolor}
\usepackage{enumitem}
\usepackage{titlesec}
\usepackage{titling}
\usepackage[hidelinks]{hyperref}

\newcolumntype{Y}{>{\raggedright\arraybackslash}X}
\renewcommand{\arraystretch}{1.35}
\setlength{\parskip}{0.6em}
\setlength{\parindent}{0pt}
\setlength{\tabcolsep}{8pt}

% Mã nguồn nội tuyến
\newcommand{\code}[1]{\texttt{#1}}
% Đường dẫn tệp: cho phép ngắt dòng tại "/" và "." để không tràn lề
\newcommand{\fpath}[1]{\nolinkurl{#1}}

% ---- Phong cách bảng -------------------------------------------------
\definecolor{navy}{RGB}{31,78,121}
\definecolor{headerblue}{RGB}{31,78,121}
\definecolor{rowalt}{RGB}{234,241,248}
\arrayrulecolor{navy!55}
% Ô tiêu đề: nền xanh (qua \rowcolor) + chữ trắng đậm
\newcommand{\thead}[1]{\textbf{\textcolor{white}{#1}}}
% ---------------------------------------------------------------------

\titleformat{\section}{\Large\bfseries\color{navy}}{\thesection.}{0.6em}{}
\titleformat{\subsection}{\large\bfseries\color{navy}}{\thesubsection}{0.6em}{}
\titlespacing*{\section}{0pt}{1.4em}{0.6em}

\hypersetup{pdftitle={Tài liệu kiểm thử NaviMart},pdfauthor={Nhóm NaviMart}}

\begin{document}

% =====================================================================
% TRANG BÌA
% =====================================================================
\begin{titlepage}
\centering
\vspace*{1.2cm}
% TODO: chỉnh tên trường/khoa cho đúng nhóm của bạn nếu cần
{\large ĐẠI HỌC BÁCH KHOA HÀ NỘI\par}
{\large Trường Công nghệ Thông tin và Truyền thông\par}
\vspace{2.4cm}
{\large\bfseries TÀI LIỆU KIỂM THỬ PHẦN MỀM\par}
{\normalsize Phiên bản 0.0.1\par}
\vspace{1.4cm}
\rule{0.8\linewidth}{0.4pt}\par
\vspace{0.6cm}
{\Huge\bfseries\color{navy} HỆ THỐNG NAVIMART\par}
\vspace{0.5cm}
{\large Quản lý kho thực phẩm gia đình, kế hoạch bữa ăn,\\[2pt] công thức và danh sách mua sắm dùng chung\par}
\vspace{0.4cm}
\rule{0.8\linewidth}{0.4pt}\par
\vspace{1.2cm}
{\large Môn: Phát triển phần mềm theo chuẩn kỹ năng ITSS\par}
\vspace{2cm}

% TODO: điền danh sách thành viên nhóm
\begin{tabular}{|l|l|}
\hline
\rowcolor{headerblue}
\thead{Họ và tên} & \thead{MSSV} \\
\hline
\rule{0pt}{2.4ex}\dots\dots\dots\dots\dots\dots & \dots\dots\dots \\
\hline
\rule{0pt}{2.4ex}\dots\dots\dots\dots\dots\dots & \dots\dots\dots \\
\hline
\rule{0pt}{2.4ex}\dots\dots\dots\dots\dots\dots & \dots\dots\dots \\
\hline
\end{tabular}

\vfill
{\itshape Hà Nội, ngày 14 tháng 6 năm 2026\par}
\end{titlepage}

\tableofcontents
\newpage

% =====================================================================
\section{Giới thiệu và mục đích tài liệu}

NaviMart là một ứng dụng quản lý gia đình xoay quanh thực phẩm: theo dõi kho thực phẩm trong nhà, nhắc hạn
sử dụng, lập kế hoạch bữa ăn, quản lý công thức nấu ăn và tạo danh sách mua sắm dùng chung giữa các thành
viên. Vì hệ thống chạm trực tiếp đến dữ liệu được nhiều người cùng chỉnh sửa --- tồn kho, quyền chia sẻ,
trạng thái thanh toán của danh sách mua sắm --- nên độ tin cậy của các quy tắc nghiệp vụ và của lớp phân
quyền là yếu tố quyết định chất lượng. Tài liệu này được biên soạn để chứng minh các quy tắc đó hoạt động
đúng, đồng thời ghi nhận một cách minh bạch những phần đã được kiểm chứng và những phần còn là rủi ro.

Tài liệu không chỉ liệt kê số lượng testcase pass/fail mà đi theo định hướng \emph{bàn giao kiểm thử}: giải
thích mục tiêu của từng loại kiểm thử, công cụ và dữ liệu sử dụng, cách đánh giá kết quả, các lỗi phát hiện
được cùng cách xử lý, và điều kiện để có thể ký duyệt bản giao. Cấu trúc trình bày tham chiếu khung tài liệu
kiểm thử theo ISO/IEC/IEEE 29119-3, nhưng được điều chỉnh cho phù hợp với một dự án thực tế gồm
\textbf{backend NestJS} và \textbf{frontend React}.

\subsection{Mục đích}

Công tác kiểm thử trong dự án hướng tới sáu mục tiêu. Thứ nhất, xác minh backend thực hiện đúng các nghiệp
vụ trọng yếu: xác thực người dùng, quản lý gia đình và chia sẻ, kho thực phẩm, công thức, kế hoạch bữa ăn,
danh sách mua sắm và báo cáo. Thứ hai, kiểm tra hợp đồng (contract) của API ở mức HTTP --- bao gồm đường
dẫn, phương thức, tham số, kiểm tra dữ liệu đầu vào và mã lỗi trả về. Thứ ba, kiểm tra logic nghiệp vụ và
các chuyển trạng thái ở mức đơn vị. Thứ tư, kiểm tra cơ chế xác thực và phân quyền (JWT, guard theo vai trò
và theo quyền gia đình). Thứ năm, kiểm tra tính bền vững của dữ liệu và các luồng nghiệp vụ xuyên suốt bằng
kiểm thử đầu-cuối trên cơ sở dữ liệu thật. Cuối cùng, đo lường độ phủ mã nguồn và đề xuất lộ trình cải tiến.

\subsection{Đối tượng sử dụng}

Tài liệu phục vụ ba nhóm người đọc với những nhu cầu khác nhau: \textbf{lập trình viên} cần biết phạm vi
test hiện có, cách chạy và bổ sung test, cũng như các lỗi cần xử lý trước khi gộp mã; \textbf{kiểm thử
viên/QA} cần đánh giá thiết kế testcase, kết quả thực thi, độ phủ và bằng chứng để ký duyệt; và
\textbf{quản lý dự án} cần cơ sở để ra quyết định về chất lượng bàn giao và ưu tiên xử lý rủi ro.

\subsection{Phạm vi}

Phạm vi của tài liệu bao gồm cả backend (mã kiểm thử trong \code{backend/src} và \code{backend/test}) và
frontend (\code{frontend/src}, \code{frontend/test}), cùng các báo cáo độ phủ sinh ra từ Jest và Vitest.
Đây là điểm mở rộng so với khung tham chiếu thông thường vốn chỉ giới hạn ở backend. Một số hạng mục nằm
ngoài phạm vi hiện tại do chưa có hạ tầng hoặc artifact: kiểm thử đầu-cuối tự động trên trình duyệt
(Playwright/Cypress), kiểm thử bảo mật động chuyên sâu (DAST/OWASP ZAP) và kiểm thử tải. Các hạng mục này
được ghi nhận trong lộ trình ở Mục~11.

% =====================================================================
\section{Định hướng ITSS và quy trình kiểm thử}

Bộ tài liệu được tổ chức bám theo các hoạt động chuẩn của một quy trình kiểm thử: lập kế hoạch, thiết kế
testcase, chuẩn bị môi trường, thực thi và thu thập bằng chứng, cuối cùng là đánh giá chất lượng cùng phân
tích rủi ro. Mỗi hoạt động được cụ thể hóa bằng các công cụ và sản phẩm thực tế của dự án, như tóm tắt
trong Bảng~\ref{tab:itss}.

\begin{table}[h]
\centering
\caption{Ánh xạ hoạt động kiểm thử sang sản phẩm thực tế của NaviMart}
\label{tab:itss}
\begin{tabularx}{\linewidth}{|l|Y|}
\hline
\rowcolor{headerblue}
\thead{Hoạt động} & \thead{Cách thể hiện trong NaviMart} \\
\hline
Lập kế hoạch & Xác định mục tiêu, phạm vi, công cụ, môi trường và tiêu chí vào/ra (Mục 3--5). \\
\hline
Thiết kế testcase & Quy ước mã testcase, ma trận truy vết theo loại kiểm thử và phụ lục danh mục chi tiết (Mục 6, Phụ lục A). \\
\hline
Tự động hóa & Jest, ts-jest, @nestjs/testing, supertest, mongodb-memory-server cho backend; Vitest, Testing Library, jsdom cho frontend. \\
\hline
Thực thi \& bằng chứng & Kết quả chạy \code{npm run test}, \code{npm run test:e2e}, \code{vitest run}; báo cáo độ phủ; nhật ký lỗi. \\
\hline
Đánh giá chất lượng & Kết luận theo từng tầng kiểm thử, sổ rủi ro và tiêu chí đóng (Mục 10--11). \\
\hline
\end{tabularx}
\end{table}

% =====================================================================
\section{Tổng quan hệ thống và phạm vi kiểm thử}

\subsection{Kiến trúc backend}

Backend NaviMart được xây dựng trên nền NestJS, cung cấp REST API kèm kênh WebSocket cho cập nhật thời gian
thực. Hệ thống áp dụng xác thực JWT không trạng thái với cặp access/refresh token, và phân quyền hai lớp:
theo vai trò (admin/user) ở cấp ứng dụng, và theo quyền của thành viên trong gia đình ở cấp tài nguyên. Dữ
liệu được lưu trên MongoDB thông qua Mongoose. Bảng~\ref{tab:modules} liệt kê các module nghiệp vụ và mức độ
được kiểm thử hiện tại; có thể thấy các module cốt lõi (xác thực, gia đình, kho, công thức, bữa ăn, danh
sách mua sắm) đều đã được phủ, trong khi một số module ngoại vi vẫn còn khoảng trống.

\begin{table}[h]
\centering
\caption{Các module backend và mức độ kiểm thử}
\label{tab:modules}
\begin{tabularx}{\linewidth}{|l|Y|c|}
\hline
\rowcolor{headerblue}
\thead{Module} & \thead{Nội dung nghiệp vụ} & \thead{Kiểm thử} \\
\hline
auth & Đăng ký, đăng nhập, refresh/rotate token, đăng xuất, quên/đặt lại mật khẩu, xác minh email. & Đầy đủ \\
\hline
families & Tạo gia đình, mã mời, gia nhập, quản lý quyền thành viên. & Đầy đủ \\
\hline
pantry & Quản lý item kho, tiêu thụ/bỏ đi, lọc theo hạn dùng, ghi sự kiện kho. & Đầy đủ \\
\hline
meals & Kế hoạch bữa ăn, gắn công thức, tính nguyên liệu thiếu, sinh danh sách mua. & Đầy đủ \\
\hline
recipes & Tìm kiếm/gợi ý theo độ phủ kho, yêu thích, duyệt nội dung, CRUD. & Đầy đủ \\
\hline
shopping-lists & Quản lý danh sách/mục, hoàn tất danh sách và đẩy vào kho. & Đầy đủ \\
\hline
reports & Báo cáo tiêu thụ, lãng phí, chi tiêu (tổng hợp). & Đầy đủ \\
\hline
catalog & Danh mục thực phẩm/nhóm/đơn vị; tìm theo tên/mã vạch. & Đầy đủ \\
\hline
notifications & Thông báo và nhắc hạn dùng. & Đầy đủ \\
\hline
users & Hồ sơ người dùng, cập nhật thông tin cá nhân. & Đầy đủ \\
\hline
admin & Quản trị danh mục/công thức/thống kê/người dùng. & Service (4 service) \\
\hline
realtime & Cổng và dịch vụ WebSocket. & Qua E2E \\
\hline
ai-chef, uploads & Gợi ý món bằng AI; tải ảnh. & Chưa có \\
\hline
\end{tabularx}
\end{table}

\subsection{Trong và ngoài phạm vi}

Trong phạm vi của đợt kiểm thử này là toàn bộ logic service, controller, hàm tiện ích, mapper, guard, khâu
kiểm tra dữ liệu và xử lý ngoại lệ ở backend; các kiểm thử đơn vị, kiểm thử API và kiểm thử đầu-cuối; cùng
các kiểm thử frontend cho hàm tiện ích, context, component và lớp gọi API. Phân quyền tại các endpoint rủi
ro cao (xác thực, gia đình, kho, danh sách mua sắm) được đặc biệt chú trọng. Ngoài phạm vi hiện tại là kiểm
thử đầu-cuối trên trình duyệt thật, kiểm thử bảo mật động chuyên sâu, kiểm thử hiệu năng/tải, và các hoạt
động vận hành sản xuất như sao lưu/phục hồi.

% =====================================================================
\section{Chiến lược kiểm thử}

Chiến lược kiểm thử được phân tầng theo mô hình tháp kiểm thử, với nguyên tắc cốt lõi: \emph{mỗi loại test
chỉ chứng minh đúng thứ nó có khả năng chứng minh}. Phân tầng rõ ràng giúp tránh ngộ nhận --- ví dụ dùng
kiểm thử đơn vị để kết luận về hành vi của cơ sở dữ liệu, hoặc dùng kiểm thử controller độc lập để khẳng
định toàn bộ chuỗi bảo mật đã đúng. Bốn tầng backend và một tầng frontend được mô tả trong
Bảng~\ref{tab:layers}.

\begin{table}[h]
\centering
\caption{Các tầng kiểm thử}
\label{tab:layers}
\begin{tabularx}{\linewidth}{|l|Y|Y|}
\hline
\rowcolor{headerblue}
\thead{Tầng} & \thead{Mục tiêu} & \thead{Công cụ \& đặc điểm} \\
\hline
Unit & Logic service, util, mapper. & Jest với model/service được mock; không khởi tạo toàn bộ Nest context. \\
\hline
API/Controller & Route, mã trạng thái, kiểm tra dữ liệu, hình dạng JSON, ánh xạ ngoại lệ. & supertest qua tiện ích \code{createApiTestApp}; service mock, guard mock. \\
\hline
Security & 401/403, quyền, guard, phân quyền cấp đối tượng. & Guard mock ở mức contract; guard \emph{thật} ở mức E2E. \\
\hline
Integration/E2E & Lưu trữ, luồng xuyên module, realtime. & App NestJS thật + mongodb-memory-server + socket.io-client. \\
\hline
Frontend & Hàm thuần, context, component, lớp API. & Vitest + Testing Library trên jsdom; mock fetch/api/socket. \\
\hline
\end{tabularx}
\end{table}

Bốn nguyên tắc được tuân thủ xuyên suốt. Một là, kiểm thử đơn vị dùng để chứng minh logic nghiệp vụ cục bộ
chứ không kết luận về giao dịch hay cơ sở dữ liệu thật. Hai là, kiểm thử API ở tầng controller chứng minh
hợp đồng HTTP với cơ chế kiểm tra dữ liệu thật, nhưng do guard và service được mock nên không tự động chứng
minh toàn bộ chuỗi bảo mật. Ba là, kiểm thử bảo mật bắt buộc phải bao gồm các đường đi tiêu cực: thiếu
token, thiếu quyền, token sai hoặc hết hạn, và truy cập tài nguyên của người khác. Bốn là, kiểm thử
đầu-cuối phải dùng nguồn dữ liệu cô lập để kết quả lặp lại được và không phụ thuộc vào cơ sở dữ liệu cục bộ.

Tương ứng với mỗi nhóm, tiêu chí vào/ra được định nghĩa rõ ràng. Nhóm Unit/API/Security yêu cầu mã biên dịch
được và spec đã định nghĩa khi bắt đầu, và yêu cầu toàn bộ test pass không còn lỗi không giải thích được khi
kết thúc. Nhóm E2E yêu cầu mongodb-memory-server khởi động được và các biến môi trường đã được thiết lập;
khi kết thúc, các luồng xuyên module phải được xác minh thành công. Cổng QA cho bản phát hành yêu cầu không
còn rủi ro mức P0 nào đang mở, và mọi rủi ro P1 đều có kế hoạch xử lý hoặc được chấp nhận có cân nhắc.

% =====================================================================
\section{Môi trường, công cụ và dữ liệu kiểm thử}

Backend được kiểm thử bằng Jest 30 (qua ts-jest 29) kết hợp tiện ích \code{@nestjs/testing}; các kiểm thử
API dùng supertest, các luồng đầu-cuối dùng mongodb-memory-server để cấp một thể hiện MongoDB cô lập ngay
trong tiến trình, và socket.io-client cho phần realtime. Frontend dùng Vitest 4.1 trên môi trường jsdom
cùng Testing Library. Độ phủ backend được thu thập bằng Jest; độ phủ frontend hiện chưa được cấu hình và
được xếp vào công việc của Giai đoạn~2. Tại thời điểm lập tài liệu, dự án \emph{chưa} có pipeline tích hợp
liên tục (CI) chạy test tự động --- đây là rủi ro được ghi nhận ở Mục~10.

Các lệnh chạy kiểm thử được chuẩn hóa như sau:

\begin{quote}\ttfamily\small
\# Backend (thư mục backend/)\\
npm run test\\
npm run test:cov\\
E2E\_USE\_MEMORY\_MONGO=true npm run test:e2e\\[4pt]
\# Frontend (thư mục frontend/)\\
npm run test
\end{quote}

Về dữ liệu kiểm thử, dự án xây dựng một bộ tiện ích dùng chung đặt ngoài thư mục \code{src} để không bị tính
vào độ phủ. Cụ thể, \code{fixtures.ts} cung cấp các hàm factory tạo dữ liệu mẫu (người dùng, gia đình, công
thức, item kho); \code{mock-model.ts} mô phỏng model Mongoose với truy vấn nối chuỗi và mặc định trả về rỗng
để lỗi ``quên mock'' lộ ra ngay; còn \code{api-test-app.ts} dựng nhanh một controller trong app Nest với
guard và pipe có thể cấu hình trạng thái xác thực/quyền theo từng test. Ở frontend, việc mock được thực hiện
trực tiếp bằng \code{vi.mock} cho lớp gọi API và socket. Riêng kiểm thử đầu-cuối dùng dữ liệu được seed thật
vào MongoDB cô lập và được dọn dẹp sau mỗi bộ test.

% =====================================================================
\section{Thiết kế testcase và ma trận truy vết}

Testcase được thiết kế dựa trên rủi ro nghiệp vụ, luồng chính, luồng lỗi, kiểm tra dữ liệu, phân quyền,
chuyển trạng thái và tính nhất quán dữ liệu. Để bảo đảm khả năng truy vết, mỗi testcase được gắn một mã có
tiền tố theo module và theo loại kiểm thử, theo quy ước trong Bảng~\ref{tab:prefix}. Các kiểm thử API hiện
có trong mã nguồn đã áp dụng sẵn quy ước này (ví dụ \code{API-AUTH-001}, \code{API-PAN-005}).

\begin{table}[h]
\centering
\caption{Quy ước mã testcase}
\label{tab:prefix}
\begin{tabularx}{\linewidth}{|l|Y|l|}
\hline
\rowcolor{headerblue}
\thead{Tiền tố} & \thead{Ý nghĩa} & \thead{Ví dụ} \\
\hline
\code{API-<MOD>} & Kiểm thử API/controller theo module. & API-FAM-003 \\
\hline
\code{UT-<MOD>} & Kiểm thử đơn vị service/util. & UT-PAN-009 \\
\hline
\code{SEC-*} & Kiểm thử bảo mật/phân quyền. & SEC-JWT-001 \\
\hline
\code{E2E-*}, \code{INT-*} & Kiểm thử đầu-cuối / tích hợp. & INT-CONC-001 \\
\hline
\code{FE-*} & Kiểm thử frontend. & FE-AUTH-001 \\
\hline
\end{tabularx}
\end{table}

Ma trận truy vết (Bảng~\ref{tab:matrix}) cho thấy mỗi module được phủ bởi những loại kiểm thử nào. Có thể
đọc ma trận này như một bản đồ độ tin cậy: các ô đầy đủ cho biết nghiệp vụ đã được chứng minh ở nhiều tầng,
trong khi các ô trống chỉ ra nơi cần ưu tiên bổ sung.

\begin{table}[h]
\centering
\small
\setlength{\tabcolsep}{5pt}
\caption{Ma trận truy vết theo loại kiểm thử}
\label{tab:matrix}
\begin{tabularx}{\linewidth}{|l|Y|Y|Y|Y|}
\hline
\rowcolor{headerblue}
\thead{Module} & \thead{Unit} & \thead{API} & \thead{Security} & \thead{Integration/E2E} \\
\hline
auth & service (21) & API (10) & 401; JWT âm tính (SEC-JWT-001..005) & đăng ký/refresh/đặt lại mật khẩu \\
\hline
families & service (13), util (4) & API (7) & 401/403 & mời/gia nhập \\
\hline
pantry & service (15), util (11) & API (8) & 401/403; object-level (SEC-OBJ-001) & lọc hạn dùng; concurrency (INT-CONC-001) \\
\hline
recipes & service (19) & API (7) & 401/403 (sai vai trò) & yêu thích/gợi ý \\
\hline
shopping-lists & service (12) & API (6) & 401/403; object-level (SEC-OBJ-002) & hoàn tất $\to$ kho \\
\hline
meals & service (24) & API (8) & 401/403; object-level (SEC-OBJ-003) & bữa ăn $\to$ thiếu $\to$ danh sách \\
\hline
reports & service (3) & API (6) & 401/403 (view\_reports) & tổng hợp báo cáo \\
\hline
catalog & service (6) & API (5) & 401 & đọc qua E2E \\
\hline
users & service (7) & API (6) & 401 & hồ sơ qua E2E \\
\hline
notifications & service (8) + nhắc hạn (6) & API (5) & 401 & một phần \\
\hline
admin & 4 service (39) & --- & promote (e2e) & nâng quyền admin (e2e) \\
\hline
\end{tabularx}
\end{table}

Về kỹ thuật thiết kế, dự án vận dụng phân hoạch tương đương và phân tích giá trị biên cho các trường số
lượng và thời hạn (ví dụ số lượng bằng/lớn hơn/nhỏ hơn tồn kho; cửa sổ nhắc hạn bằng 0 ngày), kiểm thử
chuyển trạng thái (item kho chuyển giữa các trạng thái hoạt động/đã dùng hết; danh sách mua chuyển sang hoàn
tất), kiểm thử tiêu cực (không tìm thấy, hoàn tất danh sách đã hoàn tất, sai vai trò), ma trận RBAC cho phân
quyền, và kiểm thử tương tranh cho điều kiện đua (race condition) khi tiêu thụ kho đồng thời.

% =====================================================================
\section{Kiểm thử đơn vị}

Kiểm thử đơn vị tập trung vào logic nghiệp vụ ở mức service và hàm tiện ích, với mọi phụ thuộc (model
Mongoose, service khác) đều được mock để test nhanh, cô lập và dễ truy nguyên lỗi. Đây là tầng mang lại lợi
tức cao nhất bởi nó kiểm soát trực tiếp các quy tắc nghiệp vụ và các nhánh rẽ phức tạp.

Nhóm \textbf{xác thực} được xem là ưu tiên cao nhất do liên quan trực tiếp đến danh tính và token. Hai mươi
mốt testcase bao phủ trọn vẹn vòng đời token và mật khẩu: tạo người dùng kèm gia đình sở hữu rồi phát token;
chuyển lỗi trùng khóa thành xung đột; từ chối khi sai định danh, sai mật khẩu hoặc tài khoản chưa kích hoạt;
xác minh và xoay vòng refresh token (bao gồm trường hợp sai chữ ký và hash không khớp); đặt lại mật khẩu và
xác minh email với token hợp lệ, sai và hết hạn; không tiết lộ sự tồn tại của tài khoản khi quên mật khẩu;
và cấp token tạm trong môi trường không phải sản xuất khi tắt gửi email.

Nhóm \textbf{kho thực phẩm} kết hợp mười lăm testcase cho service và mười một testcase cho hàm tiện ích tính
hạn dùng. Phần service kiểm tra việc lọc theo gia đình, từ chối tiêu thụ item không hoạt động hoặc vượt tồn,
chuyển sang trạng thái đã dùng hết khi tiêu thụ hết kèm ghi sự kiện, tái kích hoạt khi số lượng dương trở
lại, thay thế thuộc tính từ danh mục khi đổi loại thực phẩm, và đặc biệt là \emph{trừ kho nguyên tử} để
chống điều kiện đua. Phần hàm tiện ích kiểm tra việc chuẩn hóa thời gian
về nửa đêm, cộng ngày mà không làm thay đổi đầu vào, và phân loại trạng thái an toàn/sắp hết hạn/đã hết hạn
theo cửa sổ cảnh báo, bao gồm cả giá trị biên.

Các nhóm nghiệp vụ còn lại được phủ tương tự. Công thức (19 testcase) kiểm tra xếp hạng gợi ý theo tỷ lệ
khớp nguyên liệu, ưu tiên item sắp hết hạn, đếm lượt yêu thích, tự duyệt khi người tạo là admin và cấm người
dùng thường sửa công thức của người khác. Bữa ăn (24 testcase gồm cả tính nguyên liệu thiếu và sinh danh
sách mua) kiểm tra việc khớp nguyên liệu theo mã/đơn vị/tên, co giãn theo khẩu phần và bỏ qua nguyên liệu
tùy chọn. Danh sách mua sắm (12 testcase) kiểm tra việc hoàn tất danh sách và đẩy các mục đã mua vào kho, từ
chối hoàn tất danh sách đã hoàn tất hoặc không có mục nào được tích. Gia đình (13+4 testcase) kiểm tra quy
tắc phân quyền nội bộ như cấm tự sửa quyền của chính mình hay sửa quyền của chủ gia đình. Các module được bổ
sung kiểm thử trong đợt này --- danh mục, người dùng, thông báo, nhắc hạn (expiry) và bốn service quản trị
(thống kê, danh mục, công thức, người dùng) --- lấp các khoảng trống độ phủ rõ rệt nhất, đưa độ phủ tổng từ
56\% lên 68\% câu lệnh.

% =====================================================================
\section{Kiểm thử API/Controller}

Kiểm thử API xác minh hợp đồng HTTP ở tầng controller: đường dẫn, phương thức, tham số đường dẫn và truy
vấn, thân yêu cầu, khâu kiểm tra dữ liệu (loại bỏ trường lạ, ép kiểu), hình dạng phản hồi và việc ánh xạ
ngoại lệ thành mã trạng thái phù hợp. Các kiểm thử dùng tiện ích \code{createApiTestApp} để dựng controller
trong một app Nest với guard mock có thể cấu hình trạng thái xác thực và quyền theo từng kịch bản.
Bảng~\ref{tab:api} tóm tắt 68 testcase API trên mười module.

\begin{table}[h]
\centering
\caption{Danh mục kiểm thử API}
\label{tab:api}
\begin{tabularx}{\linewidth}{|l|l|Y|}
\hline
\rowcolor{headerblue}
\thead{Module} & \thead{Mã} & \thead{Kịch bản tiêu biểu} \\
\hline
auth & API-AUTH-001..010 & Đăng ký 201; thiếu mật khẩu/trường lạ 400; trùng email/sđt 409; đăng nhập 200; sai thông tin 401; đăng xuất không token 401. \\
\hline
families & API-FAM-001..007 & Lấy gia đình 200; không có gia đình 404; tạo mời thiếu quyền 403; mã mời sai 404; đã là thành viên 409; tự sửa quyền 400. \\
\hline
pantry & API-PAN-001..008 & Liệt kê 200; chưa đăng nhập 401; tạo thiếu quyền 403; số lượng âm 400; tiêu thụ quá tồn 400; không tồn tại 404. \\
\hline
recipes & API-RCP-001..007 & Tìm kiếm/gợi ý 200; công thức đã lưu trữ 404; tạo sai vai trò 403; thiếu tên 400; thêm yêu thích 201. \\
\hline
shopping-lists & API-SL-001..006 & Tạo 201; thiếu quyền 403; thiếu số lượng 400; danh sách không tồn tại 404; hoàn tất danh sách đã hoàn tất 400; hoàn tất hợp lệ 200. \\
\hline
meals & API-MEAL-001..008 & Liệt kê theo khoảng ngày 200; thiếu ngày 400; tạo thiếu quyền 403; thiếu session 400; tạo hợp lệ 201; không tồn tại 404. \\
\hline
reports & API-RPT-001..006 & Dashboard/xu hướng/lãng phí 200; chưa đăng nhập 401; thiếu quyền view\_reports 403; thiếu khoảng ngày 400. \\
\hline
catalog & API-CAT-001..005 & Danh mục/đơn vị 200; chưa đăng nhập 401; tìm món kèm query 200; limit ngoài khoảng 400. \\
\hline
users & API-USR-001..006 & Lấy hồ sơ 200; chưa đăng nhập 401; cập nhật hợp lệ 200; gender sai enum 400; trường lạ 400; không tồn tại 404. \\
\hline
notifications & API-NTF-001..005 & Liệt kê 200; chưa đăng nhập 401; limit ngoài khoảng 400; đánh dấu đã đọc/đọc tất cả 200. \\
\hline
\end{tabularx}
\end{table}

Một điểm đáng chú ý về mặt phương pháp: trong mọi kịch bản tiêu cực, kiểm thử không chỉ khẳng định mã trạng
thái trả về đúng, mà còn khẳng định service \emph{không} được gọi khi yêu cầu bị từ chối. Điều này bảo đảm
lớp kiểm tra dữ liệu và phân quyền chặn yêu cầu sai \emph{trước khi} chạm tới logic nghiệp vụ.

% =====================================================================
\section{Kết quả thực thi và độ phủ}

\subsection{Kết quả regression}

Toàn bộ bộ kiểm thử ở trạng thái xanh. Backend chạy 264 testcase đơn vị/API trên 34 bộ test cùng 32 testcase
đầu-cuối trên 4 bộ test, tất cả đều pass, không có lỗi hay thất bại nào; thời gian chạy phần đơn vị/API
khoảng mười giây. Frontend chạy 61 testcase trên 9 tệp, cũng pass toàn bộ. So với hiện trạng đầu kỳ (176
testcase đơn vị/API backend và 32 testcase frontend), bộ test đã được mở rộng đáng kể: bổ sung kiểm thử đơn
vị cho các service trước đây bỏ trống (danh mục, người dùng, thông báo, nhắc hạn, ba service quản trị), bổ
sung kiểm thử API ở mức HTTP cho năm controller (danh mục, người dùng, thông báo, bữa ăn, báo cáo), thêm
kiểm thử bảo mật và tương tranh ở backend, cùng các kiểm thử frontend mới --- đồng thời phát hiện và sửa một
lỗi nghiệp vụ thực.

\begin{table}[h]
\centering
\caption{Tổng hợp kết quả regression}
\label{tab:results}
\begin{tabularx}{\linewidth}{|Y|c|c|}
\hline
\rowcolor{headerblue}
\thead{Phạm vi} & \thead{Bộ test} & \thead{Testcase (pass/tổng)} \\
\hline
Backend đơn vị/API & 34 & 264 / 264 \\
\hline
Backend đầu-cuối & 4 & 32 / 32 \\
\hline
Frontend (Vitest) & 9 & 61 / 61 \\
\hline
\end{tabularx}
\end{table}

\subsection{Độ phủ mã nguồn backend}

Độ phủ tổng thể trên toàn bộ \code{src} đạt 68,03\% câu lệnh và 57,97\% nhánh, tăng đáng kể so với mức
56,24\% / 47,38\% đầu đợt nhờ việc lấp các module trước đây bỏ trống. Phần còn lại kéo tỷ lệ chung xuống chủ
yếu là các tệp chưa được phủ (ai-chef, uploads, cổng realtime, controller quản trị, cấu hình, seed dữ liệu)
--- nhiều trong số đó phù hợp để phủ bằng kiểm thử đầu-cuối hơn là kiểm thử đơn vị. Khi xét riêng theo module
(Bảng~\ref{tab:cov}), các module nghiệp vụ cốt lõi đều có độ phủ cao --- phản ánh đúng chiến lược ``ưu tiên
logic nghiệp vụ''.

\begin{table}[h]
\centering
\caption{Độ phủ câu lệnh theo module (backend)}
\label{tab:cov}
\begin{tabularx}{\linewidth}{|l|c|Y|}
\hline
\rowcolor{headerblue}
\thead{Module} & \thead{\% câu lệnh} & \thead{Nhận xét} \\
\hline
notifications & 88,11 & Đạt (service + API + nhắc hạn) \\
\hline
meals & 85,36 & Đạt \\
\hline
families & 84,76 & Đạt \\
\hline
users & 83,33 & Đạt (service + API) \\
\hline
pantry & 83,04 & Đạt \\
\hline
catalog & 82,14 & Đạt (service + API) \\
\hline
reports & 81,48 & Đạt (service + API) \\
\hline
recipes & 78,81 & Đạt \\
\hline
shopping-lists & 76,81 & Đạt \\
\hline
auth & 71,05 & Đạt (service cao; controller/strategy kéo xuống) \\
\hline
admin & 47,54 & Trung bình (service đã test; controller chưa) \\
\hline
realtime & 9,75 & Thấp (chỉ qua E2E) \\
\hline
ai-chef, uploads & 0,00 & Chưa có kiểm thử đơn vị \\
\hline
\end{tabularx}
\end{table}

Độ phủ frontend hiện chưa đo được do Vitest chưa được cấu hình thu thập coverage; việc bật công cụ đo được
xếp vào Giai đoạn~2.

\subsection{Phân tích lý do}

Phân bố độ phủ phản ánh ba lựa chọn có chủ đích. Thứ nhất, kiểm thử đơn vị được dồn vào tầng service --- nơi
chứa hầu hết quy tắc nghiệp vụ và nhánh rẽ --- nên các module cốt lõi đạt độ phủ cao, khóa chặt rủi ro sai
lệch nghiệp vụ. Thứ hai, việc cố ép độ phủ cho các tệp cấu hình, seed, cổng realtime hay DTO mang lại lợi
tức thấp về phát hiện lỗi và làm tăng chi phí bảo trì, nên những phần này được ưu tiên phủ bằng kiểm thử
đầu-cuối thay vì kiểm thử đơn vị. Thứ ba, các khoảng trống còn lại (admin, uploads, ai-chef) được xác định
rõ là hạng mục ưu tiên khi thiết lập cổng chất lượng tăng dần trong các giai đoạn sau.

% =====================================================================
\section{Rủi ro, lỗi và tiêu chí đóng}

\subsection{Sổ rủi ro}

Sau đợt kiểm thử, hai rủi ro mức P0 quan trọng nhất đã được đóng: lỗi tương tranh trừ kho và sự thiếu hụt
kiểm thử bảo mật âm tính. Các rủi ro còn lại chủ yếu thuộc về hạ tầng kiểm thử --- thiếu CI, thiếu cổng độ
phủ, frontend chưa đo coverage --- và được quy hoạch vào Giai đoạn~2 (Bảng~\ref{tab:risk}).

\begin{table}[h]
\centering
\caption{Sổ rủi ro kiểm thử}
\label{tab:risk}
\begin{tabularx}{\linewidth}{|l|l|Y|l|}
\hline
\rowcolor{headerblue}
\thead{ID} & \thead{Mức} & \thead{Mô tả / ảnh hưởng} & \thead{Trạng thái} \\
\hline
TEST-004 & High & Lỗi tương tranh trừ kho gây oversell. & \textbf{Đã đóng} \\
\hline
TEST-005 & Med & Thiếu kiểm thử JWT âm tính và phân quyền cấp đối tượng. & \textbf{Đã đóng} \\
\hline
TEST-003 & Med & Còn vài phần chưa có kiểm thử (ai-chef, uploads, cổng realtime, controller quản trị). & Giảm thiểu phần lớn \\
\hline
TEST-001 & High & Chưa có CI chạy test $\Rightarrow$ PR lỗi vẫn có thể gộp. & Mở (Giai đoạn 2) \\
\hline
TEST-002 & Med & Frontend chưa cấu hình độ phủ. & Mở (Giai đoạn 2) \\
\hline
TEST-006 & Med & Chưa có cổng chất lượng độ phủ. & Mở (Giai đoạn 2) \\
\hline
TEST-007 & Low & Mock frontend gắn chặt cài đặt (chưa dùng MSW). & Theo dõi \\
\hline
\end{tabularx}
\end{table}

\subsection{Lỗi đã ghi nhận}

Đợt kiểm thử ghi nhận một lỗi nghiệp vụ thực và đã khắc phục: \textbf{BUG-CONC-001} --- hàm tiêu thụ kho theo
mẫu đọc--sửa--ghi gây trừ vượt khi tiêu thụ đồng thời, đã sửa bằng cập nhật nguyên tử có điều kiện và được
bảo vệ bởi INT-CONC-001 cùng một kiểm thử đơn vị tương ứng. Hai mục còn lại liên quan đến hạ tầng (thiếu
pipeline CI và thiếu cổng độ phủ) được chuyển sang Giai đoạn~2.

\subsection{Tiêu chí đóng P0}

Hai tiêu chí P0 liên quan đến chất lượng nghiệp vụ và bảo mật đã được thỏa mãn: kiểm thử tương tranh chứng
minh hệ thống chống trừ vượt, và kiểm thử bảo mật âm tính (JWT, phân quyền cấp đối tượng) đã có bằng chứng.
Hai tiêu chí còn lại --- thiết lập CI chạy test cho cả backend lẫn frontend, và thiết lập cổng độ phủ chống
suy giảm --- thuộc về Giai đoạn~2.

% =====================================================================
\section{Đánh giá mức độ đáp ứng và kết luận}

\subsection{Đánh giá theo từng tầng}

Xét theo từng tầng kiểm thử, có thể đánh giá như sau. Kiểm thử đơn vị ở mức \emph{tốt}: các service cốt lõi
được phủ cao với bộ fixture và mock chuẩn hóa. Kiểm thử API cũng ở mức \emph{tốt}: hợp đồng được kiểm chứng
đầy đủ cho năm module chính và có mã truy vết rõ ràng. Kiểm thử bảo mật đạt mức \emph{tốt} sau khi bổ sung
các đường đi âm tính qua guard thật. Kiểm thử tích hợp/đầu-cuối ở mức \emph{tốt} với các luồng xuyên module,
realtime và tương tranh chạy trên cơ sở dữ liệu cô lập. Riêng hạng mục độ phủ và CI \emph{cần cải thiện} và
là trọng tâm của Giai đoạn~2.

\subsection{Kết luận}

Bộ kiểm thử NaviMart đã có nền tảng vững cho regression ở tầng service và controller, đồng thời sở hữu một
lớp kiểm thử đầu-cuối thực sự với MongoDB cô lập và WebSocket. Toàn bộ kiểm thử ở trạng thái xanh và quá
trình kiểm thử đã chứng tỏ giá trị thực tiễn khi phát hiện rồi khắc phục một lỗi tương tranh nghiêm trọng.
Tuy vậy, do phần hạ tầng kiểm thử (tích hợp liên tục, cổng độ phủ và đo độ phủ frontend) chưa hoàn tất, hệ
thống chưa nên được coi là đạt QA toàn diện cho bản phát hành; những hạng mục này được chủ động hoãn sang
Giai đoạn~2 theo quyết định của nhóm.

\subsection{Lộ trình cải tiến}

Bảng~\ref{tab:roadmap} tổng hợp lộ trình. Các hạng mục Giai đoạn~1 đã hoàn thành trong đợt này; các hạng mục
Giai đoạn~2 (hạ tầng) được hoãn lại theo kế hoạch.

\begin{table}[h]
\centering
\caption{Lộ trình cải tiến kiểm thử}
\label{tab:roadmap}
\begin{tabularx}{\linewidth}{|c|l|Y|}
\hline
\rowcolor{headerblue}
\thead{Giai đoạn} & \thead{Ưu tiên} & \thead{Hạng mục} \\
\hline
1 & P0 & Kiểm thử tương tranh và sửa lỗi oversell --- \emph{đã làm}. \\
\hline
1 & P1 & Kiểm thử JWT âm tính và phân quyền cấp đối tượng qua guard thật --- \emph{đã làm}. \\
\hline
1 & P1 & Bổ sung kiểm thử cho danh mục, người dùng, thông báo, thống kê quản trị --- \emph{đã làm}. \\
\hline
1 & P2 & Mở rộng kiểm thử frontend (lớp API, trang chính) --- \emph{đã làm}. \\
\hline
2 & P0 & Thiết lập pipeline CI chạy test backend và frontend. \\
\hline
2 & P0 & Bật đo độ phủ frontend và đặt cổng độ phủ cho cả hai phía. \\
\hline
--- & P2 & Bổ sung E2E trình duyệt (Playwright) và quét DAST (OWASP ZAP). \\
\hline
\end{tabularx}
\end{table}

% =====================================================================
\appendix
\section{Phụ lục A. Danh mục testcase chi tiết}

\subsection{Backend --- kiểm thử đơn vị}
\begin{longtable}{|>{\raggedright\arraybackslash}p{0.40\linewidth}|>{\raggedright\arraybackslash}p{0.50\linewidth}|}
\hline
\rowcolor{headerblue}
\thead{Tệp} & \thead{Nội dung (số testcase)} \\
\hline
\endhead
\fpath{auth.service.spec.ts} & Vòng đời token và mật khẩu: tạo người dùng + gia đình, phát/xoay token, đặt lại và xác minh email, không lộ tồn tại tài khoản (21). \\
\hline
\fpath{families.service.spec.ts} & Tạo gia đình, gia nhập theo mã mời, các quy tắc phân quyền nội bộ (13). \\
\hline
\fpath{family-access.util.spec.ts} & Từ chối khi không hoạt động/không phải thành viên/đã rời (4). \\
\hline
\fpath{pantry.service.spec.ts} & CRUD kho, tiêu thụ/bỏ đi, sự kiện kho, và trừ kho nguyên tử chống tương tranh (15). \\
\hline
\fpath{expiry-status.util.spec.ts} & Tính và phân loại hạn dùng, gồm giá trị biên (11). \\
\hline
\fpath{meals.service.spec.ts} & Kế hoạch bữa ăn, gắn công thức, ủy quyền tính nguyên liệu thiếu (11). \\
\hline
\fpath{missing-ingredients.service.spec.ts} & Khớp nguyên liệu, co giãn theo khẩu phần, xử lý nguyên liệu tùy chọn (10). \\
\hline
\fpath{shopping-list-generation.service.spec.ts} & Sinh danh sách từ nguyên liệu thiếu (4). \\
\hline
\fpath{recipes.service.spec.ts} & Xếp hạng gợi ý, yêu thích, duyệt nội dung, phân quyền chỉnh sửa (19). \\
\hline
\fpath{shopping-lists.service.spec.ts} & Quản lý danh sách, hoàn tất và đẩy vào kho (12). \\
\hline
\fpath{reports.service.spec.ts} & Tổng hợp tiêu thụ/lãng phí/bảng điều khiển (3). \\
\hline
\fpath{catalog.service.spec.ts} & Danh mục, tìm theo tên/mã vạch, ánh xạ phản hồi (6). \\
\hline
\fpath{users.service.spec.ts} & Hồ sơ, cập nhật thông tin và thiết lập thông báo (7). \\
\hline
\fpath{notifications.service.spec.ts} & Liệt kê, đánh dấu đã đọc, tạo có khử trùng lặp (8). \\
\hline
\fpath{expiry-notifications.service.spec.ts} & Quét item hết hạn, phân loại expired/expiring/safe, dedupe, phát realtime (6). \\
\hline
\fpath{admin-stats.service.spec.ts} & Tổng hợp thống kê người dùng/gia đình/công thức (2). \\
\hline
\fpath{admin-recipes.service.spec.ts} & Phân trang kiểm duyệt, cập nhật status, chặn công thức đã lưu trữ (8). \\
\hline
\fpath{admin-users.service.spec.ts} & Lọc/tìm kiếm, tạo (băm mật khẩu, Conflict), cập nhật, xóa mềm (12). \\
\hline
\fpath{admin-catalog.service.spec.ts} & CRUD danh mục/món/đơn vị: slugify, kiểm tra category, Conflict, xóa mềm (13). \\
\hline
\end{longtable}

\subsection{Backend --- kiểm thử API, đầu-cuối và controller}
Controller (mock service): \code{auth} (4), \code{families} (3), \code{pantry} (4), \code{recipes} (4). \\
API (supertest, 68): API-AUTH-001..010, API-FAM-001..007, API-PAN-001..008, API-RCP-001..007,
API-SL-001..006, API-MEAL-001..008, API-RPT-001..006, API-CAT-001..005, API-USR-001..006, API-NTF-001..005. \\
Đầu-cuối: \code{main-flows} (\textasciitilde 10 luồng), \code{extended-flows} (\textasciitilde 14 luồng),
\code{security} (SEC-JWT-001..005 và SEC-OBJ-000..003, 9 test), \code{concurrency} (INT-CONC-001, 1 test).

\subsection{Frontend (Vitest)}
\begin{longtable}{|>{\raggedright\arraybackslash}p{0.42\linewidth}|>{\raggedright\arraybackslash}p{0.48\linewidth}|}
\hline
\rowcolor{headerblue}
\thead{Tệp} & \thead{Nội dung (số testcase)} \\
\hline
\endhead
\fpath{utils/expiry.test.ts} & Tính số ngày tới hạn/còn lại/quá hạn (8). \\
\hline
\fpath{api/client.test.ts} & Lưu/đọc token, làm mới khi 401, phát sự kiện chưa xác thực (9). \\
\hline
\fpath{components/ProtectedRoute.test.tsx} & Điều hướng theo trạng thái đăng nhập và vai trò (4). \\
\hline
\fpath{components/FoodAutocomplete.test.tsx} & Gợi ý có trễ và chọn món (2). \\
\hline
\fpath{contexts/AuthContext.test.tsx} & Đăng nhập/đăng ký/đăng xuất, làm mới phiên (6). \\
\hline
\fpath{contexts/DialogContext.test.tsx} & Hộp thoại cảnh báo và xác nhận (5). \\
\hline
\fpath{api/index.test.ts} & Lớp gọi API (auth/pantry/recipes/shopping-lists/catalog/users/family/notifications) ánh xạ đúng đường dẫn và tùy chọn (25). \\
\hline
\fpath{components/Button.test.tsx} & Render children, bắt onClick, vô hiệu khi disabled, hiển thị icon (4). \\
\hline
\fpath{pages/Login.test.tsx} & Kiểm tra dữ liệu, điều hướng theo vai trò, hiển thị lỗi (4). \\
\hline
\end{longtable}

% =====================================================================
\section{Phụ lục B. Checklist chạy và nghiệm thu}
\begin{tabularx}{\linewidth}{|c|l|Y|}
\hline
\rowcolor{headerblue}
\thead{Bước} & \thead{Lệnh / hoạt động} & \thead{Kết quả mong đợi} \\
\hline
1 & \code{cd backend \&\& npm run test} & 34 bộ / 264 testcase pass. \\
\hline
2 & \code{cd backend \&\& npm run test:cov} & Sinh báo cáo độ phủ trong \code{backend/coverage}. \\
\hline
3 & \code{E2E\_USE\_MEMORY\_MONGO=true npm run test:e2e} & 4 bộ / 32 testcase đầu-cuối pass. \\
\hline
4 & \code{cd frontend \&\& npm run test} & 9 tệp / 61 testcase pass. \\
\hline
5 & (Giai đoạn 2) Kiểm tra pipeline CI & Job backend và frontend chạy test/độ phủ thành công. \\
\hline
6 & Rà soát sổ rủi ro & Không còn rủi ro P0 đang mở. \\
\hline
\end{tabularx}

\end{document}
